\section{Methodological Framework}

The project follows an implementation-and-evaluation methodology. First, the university threat model is defined: students may access public study material, while protected exam artifacts are restricted to teachers. Second, the system architecture is implemented with separate ingestion, retrieval, privacy, Bloom classification, and local generation modules. Third, each component is evaluated with measurable metrics. Finally, the results are interpreted conservatively according to the available evidence.

The methodology has five stages:
\begin{enumerate}
    \item Problem formulation and requirement analysis.
    \item Data ingestion and metadata preservation.
    \item Model and retrieval pipeline development.
    \item Privacy, QA/RAG, OCR, and Bloom evaluation.
    \item Report preparation and evidence-aligned claim formulation.
\end{enumerate}

\section{Tools, Libraries, and Technologies Used}

\begin{table}[H]
\centering
\caption{Tools and technologies used in the project}
\label{tab:tools}
\begin{tabular}{p{0.30\textwidth}p{0.58\textwidth}}
\toprule

\textbf{Tool/Technology} & \textbf{Purpose}\\
\midrule
Python project modules & Implementation of ingestion, retrieval, evaluation, and app logic.\\ \hline
FAISS & Local vector index storage and retrieval.\\ \hline
PyMuPDF & PDF text extraction.\\ \hline
pytesseract/Tesseract & OCR-backed image ingestion.\\ \hline
llama.cpp CLI & CPU-only inference for \qwenmodel{}.gguf.\\ \hline
Qwen2.5-1.5B-Instruct-Q4\_K\_M.gguf & Local quantized tiny LLM used for generation.\\ \hline
Streamlit & Interactive demo interface.\\ \hline
JSON/CSV result files & Reproducible evaluation output.\\ \hline
LaTeX/Overleaf & Report preparation and academic formatting.\\ 
\bottomrule
\end{tabular}
\end{table}

\section{Data Collection / Dataset Description}

The project uses multiple evaluation artifacts:
\begin{itemize}
    \item \textbf{Figshare Bloom dataset:} primary Bloom classification dataset (2522 raw rows; 2330 unique questions after auditing). Used for LoRA training, official test evaluation, and 15\% hold-out baselines (\texttt{figshare\_combined\_dataset.csv}).
    \item \textbf{Protected exam micro-corpus:} synthetic protected artifacts used to evaluate student leakage attacks, benign student prompts, and teacher moderation prompts.
    \item \textbf{Offline academic QA benchmark:} ten local academic questions used for retrieval and answer-grounding evaluation.
    \item \textbf{PDF/image smoke set:} two multimodal cases used to test PDF and OCR-backed image RAG.
    \item \textbf{Synthetic typed images:} used to evaluate OCR token-F1, word error rate, modality preservation, and access-level preservation.
\end{itemize}

\section{Preprocessing Steps}

The Bloom dataset preprocessing normalizes label names from older Bloom categories to modern labels. For example, Knowledge maps to Remember, Comprehension maps to Understand, Application maps to Apply, Analysis maps to Analyze, Evaluation maps to Evaluate, and Synthesis maps to Create. Exact duplicate questions are removed. If any question text maps to multiple Bloom labels, conflicting duplicates are removed. The final Figshare dataset contains 2330 unique questions.

For document ingestion, text and Markdown files are chunked directly. PDFs are parsed by page, and image files are converted to text using OCR. Metadata is attached to every chunk so that modality and access level remain available during retrieval and screening.

\section{Model Development / Algorithm Design}

\subsection{Bloom Evaluation Protocol}
Bloom models are compared using \texttt{bloom\_evaluation.py} and \texttt{build\_bloom\_comparison.py}. A 15\% stratified hold-out ($n{=}379$) evaluates TF-IDF + LinearSVC and zero-shot Qwen GGUF (\texttt{bloom\_prompt.zero\_shot\_bloom\_label}). The deployed Qwen2.5-1.5B LoRA classifier is trained with \texttt{train\_qwen\_bloom.py} and evaluated on the official audited test split ($n{=}2330$, \texttt{evaluation\_results/metrics.json}). Labels come from the LoRA head; teacher reason/rewrite text is generated by GGUF in \texttt{bloom\_prompt.py}.

\subsection{RAG Pipeline}
The RAG pipeline retrieves relevant chunks from the public vector store for student queries. The local LLM then generates an answer from retrieved context. If the retrieved context does not support the answer, the system is expected to decline or avoid unsupported claims. The Qwen GGUF evaluation uses CPU-only llama.cpp inference with context size 512, maximum output length 32 tokens, two threads, no GPU layers, and deterministic decoding.

\subsection{Privacy Guard}
The privacy guard implements role-aware access control and leakage screening. Student prompts are checked for reconstruction intent and protected artifact references. Generated outputs are checked for copied spans, n-gram containment, and semantic concept overlap. Teacher moderation prompts are allowed because the teacher role is authorized to inspect protected materials.

\section{System Implementation / Modules Description}

\subsection{Ingestion Module}
The ingestion module supports text, PDF, and OCR-backed image files. It produces structured chunks with metadata. This module is responsible for ensuring that access level and modality are not lost.

\subsection{Retrieval Module}
The retrieval module stores and retrieves document chunks using local vector indexes. It supports persistent save/load behavior so public and protected vector databases can be reused.

\subsection{Multimodal RAG Module}
The multimodal RAG module wraps document ingestion, retrieval, and guarded answer generation for PDF and image inputs. PDF RAG uses extracted text, while image RAG uses OCR-derived text.

\subsection{Qwen GGUF Module}
The Qwen module runs \qwenmodel{}.gguf through standalone llama.cpp CLI. This avoids reliance on Python llama-cpp bindings that require a native compiler toolchain. The model file size is 986.049 MB.

\subsection{Streamlit Application}
The Streamlit application provides an interactive user interface for uploading files into either the public learning corpus or protected exam corpus. It preserves metadata and routes user interactions according to role.

\section{Hardware and Software Requirements}

\subsection{Development Environment}
The project is developed as a Python-based local software system. Required software includes Python dependencies for document ingestion, FAISS vector retrieval, OCR, Streamlit, and llama.cpp CLI for GGUF model inference.

\subsection{Deployment Environment}
The intended deployment requires a CPU-capable machine with sufficient memory to load a sub-1GB quantized GGUF model plus runtime context buffers. An OCR backend such as Tesseract is required for image ingestion. No custom hardware is required.

\section{Summary}

This chapter described the methodology, datasets, preprocessing, model development, implementation modules, and hardware/software requirements. The next chapter presents measured results for Bloom classification, privacy-guard behavior, QA/RAG grounding, and OCR-backed multimodal ingestion.
